% ============================================================
% 片段文件：第 1 章  集合与常用逻辑用语
% 由主文件 讲义.tex 通过 \input 调用
% 注意：本文件不含 \chapter，章标题写在主文件里
% ============================================================

\section{集合的概念}
\subsection{元素与集合}
\begin{definition}[元素]
元素是构成集合的对象。
\begin{example}
    \begin{itemize}
        \item 数字 1、2、3 是集合 $\{1,2,3\}$ 的元素。
        \item 字母 a、b、c 是集合 $\{a,b,c\}$ 的元素。
    \end{itemize}
\end{example}
\end{definition}
\begin{axiom}[集合的确定性]
给定一个集合，任何一个对象要么是这个集合的元素，要么不是这个集合的元素，二者必居其一。集合与其元素之间的关系由 \textbf{ZF} 公理系统刻画（如外延公理：两个集合相等当且仅当它们含有完全相同的元素）。
\end{axiom}
\begin{definition}[集合]
集合是具有某种共同特征的对象的全体。
\end{definition}

\begin{definition}[元素与集合的关系]
若 $a$ 是集合 $A$ 的元素，记作 $a \in A$；若 $a$ 不是集合 $A$ 的元素，记作 $a \notin A$。
\end{definition}
\begin{definition}[常用数集]
\begin{itemize}
    \item 全体非负整数组成的集合称为自然数集，记作 $\mathbb{N}$；
    \item 全体正整数组成的集合称为正整数集，记作 $\mathbb{N}^*$ 或 $\mathbb{N}_+$；
    \item 全体整数组成的集合称为整数集，记作 $\mathbb{Z}$；
    \item 全体有理数组成的集合称为有理数集，记作 $\mathbb{Q}$；
    \item 全体实数组成的集合称为实数集，记作 $\mathbb{R}$。
\end{itemize}
\end{definition}
\subsection{集合的表示方法}
\begin{definition}[列举法]
把集合中的所有元素一一列举出来，并用花括号 $\{\ \}$ 括起来表示集合的方法叫做列举法。
\end{definition}
\begin{definition}[描述法]
用集合所含元素的共同特征表示集合的方法叫做描述法，其一般形式为 $\{x \mid p(x)\}$，其中 $x$ 为元素，$p(x)$ 为元素满足的条件。
\end{definition}
\begin{definition}[图示法]
用平面上封闭曲线的内部表示集合的方法叫做图示法（韦恩图）。
\end{definition}
\subsection{集合中元素的特性}
\begin{definition}[确定性、互异性、无序性]
集合中的元素具有如下三个特性：
\begin{itemize}
    \item 确定性：给定集合，任何一个对象是不是它的元素是确定的；
    \item 互异性：集合中的元素互不相同；
    \item 无序性：集合中的元素没有先后顺序。
\end{itemize}
\end{definition}

\section{集合间的基本关系}
\subsection{子集与真子集}
\begin{definition}[子集]
若集合 $A$ 中的任意一个元素都是集合 $B$ 的元素，则称 $A$ 是 $B$ 的子集，记作 $A \subseteq B$（或 $B \supseteq A$）。\\
规定：空集是任何集合的子集。
\end{definition}
\begin{theorem}[子集的性质]
对于任意集合 $A$、$B$、$C$：
\begin{itemize}
    \item 自反性：$A \subseteq A$；
    \item 反对称性：若 $A \subseteq B$ 且 $B \subseteq A$，则 $A=B$；
    \item 传递性：若 $A \subseteq B$ 且 $B \subseteq C$，则 $A \subseteq C$。
\end{itemize}
\end{theorem}
\begin{corollary}[子集个数]
含有 $n$ 个元素的集合共有 $2^n$ 个子集，其中真子集有 $2^n-1$ 个，非空真子集有 $2^n-2$ 个。
\end{corollary}
\begin{definition}[真子集]
若 $A \subseteq B$，且存在 $b \in B$ 但 $b \notin A$，则称 $A$ 是 $B$ 的真子集，记作 $A \subsetneq B$（或 $B \supsetneq A$）。
\end{definition}
\subsection{集合相等}
\begin{definition}[集合相等]
若 $A \subseteq B$ 且 $B \subseteq A$，则称集合 $A$ 与集合 $B$ 相等，记作 $A = B$。
\end{definition}
\subsection{空集}
\begin{definition}[空集]
不含任何元素的集合叫做空集，记作 $\varnothing$。
\end{definition}

\section{集合的基本运算}
\subsection{并集}
\begin{definition}[并集]
由所有属于集合 $A$ 或属于集合 $B$ 的元素组成的集合，叫做 $A$ 与 $B$ 的并集，记作 $A \cup B$，即 $A \cup B = \{x \mid x \in A \text{ 或 } x \in B\}$。
\end{definition}
\begin{lemma}[并集的性质]
对于任意集合 $A$、$B$：
\begin{itemize}
    \item 交换律：$A \cup B=B \cup A$；
    \item 幂等律：$A \cup A=A$；
    \item $A \cup \varnothing=A$，$A \subseteq A \cup B$。
\end{itemize}
\end{lemma}
\subsection{交集}
\begin{definition}[交集]
由所有既属于集合 $A$ 又属于集合 $B$ 的元素组成的集合，叫做 $A$ 与 $B$ 的交集，记作 $A \cap B$，即 $A \cap B = \{x \mid x \in A \text{ 且 } x \in B\}$。
\end{definition}
\begin{theorem}[并集与交集的运算律]
对于任意集合 $A$、$B$、$C$：
\begin{itemize}
    \item 交换律：$A \cap B=B \cap A$；
    \item 结合律：$(A \cup B) \cup C=A \cup (B \cup C)$，$(A \cap B) \cap C=A \cap (B \cap C)$；
    \item 分配律：$A \cap (B \cup C)=(A \cap B) \cup (A \cap C)$，$A \cup (B \cap C)=(A \cup B) \cap (A \cup C)$。
\end{itemize}
\end{theorem}
\subsection{补集}
\begin{definition}[全集与补集]
若集合 $U$ 含有我们所研究问题中涉及的所有元素，则称 $U$ 为全集；对于集合 $A \subseteq U$，由 $U$ 中不属于 $A$ 的所有元素组成的集合，叫做 $A$ 在 $U$ 中的补集，记作 $\complement_U A$，即 $\complement_U A = \{x \mid x \in U \text{ 且 } x \notin A\}$。
\end{definition}
\begin{theorem}[补集的性质（德摩根律）]
设全集为 $U$，$A$、$B \subseteq U$：
\begin{itemize}
    \item $A \cup \complement_U A=U$，$A \cap \complement_U A=\varnothing$；
    \item $\complement_U(\complement_U A)=A$；
    \item $\complement_U(A \cup B)=\complement_U A \cap \complement_U B$，$\complement_U(A \cap B)=\complement_U A \cup \complement_U B$。
\end{itemize}
\end{theorem}

\section{充分条件与必要条件}
\subsection{命题及其真假}
\begin{definition}[命题]
可以判断真假的陈述句叫做命题。判断为真的命题叫做真命题，判断为假的命题叫做假命题。
\end{definition}
\subsection{充分条件、必要条件}
\begin{definition}[充分条件与必要条件]
对于命题“若 $p$，则 $q$”，若该命题为真命题，即 $p \Rightarrow q$，则称 $p$ 是 $q$ 的充分条件，$q$ 是 $p$ 的必要条件。
\end{definition}
\subsection{充要条件}
\begin{definition}[充要条件]
若 $p \Rightarrow q$ 且 $q \Rightarrow p$，即 $p \Leftrightarrow q$，则称 $p$ 是 $q$ 的充分必要条件，简称充要条件。
\end{definition}
\begin{definition}[充分不必要、必要不充分、既不充分也不必要]
\begin{itemize}
    \item 若 $p \Rightarrow q$ 但 $q \nRightarrow p$，则 $p$ 是 $q$ 的充分不必要条件；
    \item 若 $q \Rightarrow p$ 但 $p \nRightarrow q$，则 $p$ 是 $q$ 的必要不充分条件；
    \item 若 $p \nRightarrow q$ 且 $q \nRightarrow p$，则 $p$ 是 $q$ 的既不充分也不必要条件。
\end{itemize}
\end{definition}

\section{全称量词与存在量词}
\subsection{全称量词与全称量词命题}
\begin{definition}[全称量词与全称量词命题]
“所有的”“任意一个”等短语叫做全称量词，用符号“$\forall$”表示；含有全称量词的命题叫做全称量词命题，一般形式为“$\forall x \in M,\ p(x)$”。
\end{definition}
\subsection{存在量词与存在量词命题}
\begin{definition}[存在量词与存在量词命题]
“存在一个”“至少有一个”等短语叫做存在量词，用符号“$\exists$”表示；含有存在量词的命题叫做存在量词命题，一般形式为“$\exists x \in M,\ p(x)$”。
\end{definition}
\subsection{全称量词命题与存在量词命题的否定}
\begin{definition}[全称量词命题的否定]
全称量词命题“$\forall x \in M,\ p(x)$”的否定是存在量词命题“$\exists x \in M,\ \neg p(x)$”。
\end{definition}
\begin{definition}[存在量词命题的否定]
存在量词命题“$\exists x \in M,\ p(x)$”的否定是全称量词命题“$\forall x \in M,\ \neg p(x)$”。
\end{definition}
\begin{theorem}[命题与其否定]
对任意命题 $p$，命题 $p$ 与其否定 $\neg p$ 必一真一假，不能同时为真也不能同时为假。
\end{theorem}
\begin{corollary}[全称量词命题为假的判定]
全称量词命题“$\forall x \in M,\ p(x)$”为假，当且仅当存在 $x_0 \in M$ 使 $p(x_0)$ 为假。
\end{corollary}
